% =====================================================================
% Round 2: Mathematical Improvements for "The Geometry of Linear Regression"
% =====================================================================
%
% This file contains complete, self-contained LaTeX blocks to be inserted
% into lecture_notes.tex. Each block is preceded by an insertion directive.
% All blocks use the macros already defined in the main document:
%   \bX, \by, \bH, \bM, \be, \yhat, \bhat, \btilde,
%   \bI, \bQ, \bD, \Col, \ip, \norm, \R, \tr, \rank, \diag, \argmin
%
% =====================================================================


% =====================================================================
% ITEM 1: Proof of the Projection Theorem (finite-dimensional)
% INSERT: Immediately after Theorem 1.1 (the statement of the Projection
%         Theorem), replacing the current gap before "Ordinary least
%         squares is the Projection Theorem applied to..."
% =====================================================================

\begin{proof}
We work in $V = \R^n$ with the Euclidean inner product, and $S$ is a
$p$-dimensional subspace (hence closed).  The proof has three parts:
existence, uniqueness, and the orthogonality characterisation.

\medskip
\noindent\textbf{Step 1: Existence.}
Let $d = \inf_{\bm{s}\in S}\norm{\by - \bm{s}}$.
Choose any $\bm{s}_0 \in S$. Then every minimising candidate satisfies
$\norm{\by - \bm{s}} \le \norm{\by - \bm{s}_0}$,
so the search is confined to the closed, bounded set
$\{\bm{s}\in S : \norm{\by - \bm{s}} \le \norm{\by - \bm{s}_0}\}$.
Since $S$ is a subspace of a finite-dimensional space, this set is compact.
The continuous function $\bm{s}\mapsto\norm{\by - \bm{s}}$ attains its
minimum on a compact set, so there exists $\yhat\in S$ with
$\norm{\by - \yhat} = d$.

\medskip
\noindent\textbf{Step 2: Uniqueness.}
Suppose $\yhat$ and $\yhat'$ are both minimisers.
By the parallelogram law,
\[
  \norm{(\by - \yhat) + (\by - \yhat')}^2
  + \norm{(\by - \yhat) - (\by - \yhat')}^2
  = 2\norm{\by - \yhat}^2 + 2\norm{\by - \yhat'}^2.
\]
The left-hand side rewrites as
$\norm{2\by - \yhat - \yhat'}^2 + \norm{\yhat' - \yhat}^2$,
and the right-hand side equals $4d^2$.
Now $\frac{1}{2}(\yhat + \yhat') \in S$ (since $S$ is a subspace), so
\[
  \norm{2\by - \yhat - \yhat'}^2
  = 4\,\norm{\by - \tfrac{\yhat+\yhat'}{2}}^2
  \ge 4d^2.
\]
Therefore $\norm{\yhat' - \yhat}^2 \le 4d^2 - 4d^2 = 0$,
which forces $\yhat' = \yhat$.

\medskip
\noindent\textbf{Step 3: Orthogonality condition.}
Let $\bm{v}\in S$ be arbitrary and $t\in\R$.
Since $\yhat + t\bm{v}\in S$, the minimality of $\yhat$ gives
\[
  f(t)
  \;=\; \norm{\by - \yhat - t\bm{v}}^2
  \;=\; \norm{\by - \yhat}^2
        - 2t\,\ip{\by - \yhat}{\bm{v}}
        + t^2\norm{\bm{v}}^2
  \;\ge\; f(0)
  \quad\text{for all } t\in\R.
\]
The function $f$ is a quadratic in $t$ with minimum at $t=0$.
Setting $f'(0) = 0$:
\[
  f'(0) = -2\,\ip{\by-\yhat}{\bm{v}} = 0
  \qquad\Longrightarrow\qquad
  \ip{\by - \yhat}{\bm{v}} = 0.
\]
Since $\bm{v}\in S$ was arbitrary, $(\by - \yhat)\perp S$.

\medskip
\noindent\textbf{Conversely}, if $\yhat\in S$ satisfies $(\by-\yhat)\perp S$,
then for any $\bm{s}\in S$:
\[
  \norm{\by - \bm{s}}^2
  = \norm{(\by - \yhat) + (\yhat - \bm{s})}^2
  = \norm{\by - \yhat}^2 + \norm{\yhat - \bm{s}}^2
  \ge \norm{\by - \yhat}^2,
\]
where the cross term vanishes because $(\by - \yhat)\perp(\yhat-\bm{s})$
(since $\yhat - \bm{s}\in S$), and the Pythagorean theorem applies.
Hence $\yhat$ is the unique minimiser.
\end{proof}


% =====================================================================
% ITEM 2: Full FWL Proof (block elimination)
% INSERT: Replace the current proof of Theorem (FWL) in Section 5
% =====================================================================

\begin{proof}
The normal equations $\bX^\top\bX\,\bhat = \bX^\top\by$ in block form are
\begin{equation}\label{eq:fwl-blocks}
  \begin{pmatrix}
    \bX_1^\top\bX_1 & \bX_1^\top\bX_2 \\
    \bX_2^\top\bX_1 & \bX_2^\top\bX_2
  \end{pmatrix}
  \begin{pmatrix}\bhat_1 \\ \bhat_2\end{pmatrix}
  =
  \begin{pmatrix}\bX_1^\top\by \\ \bX_2^\top\by\end{pmatrix}.
\end{equation}

\noindent\textbf{Step 1: Solve the first block for $\bhat_1$.}
The first block equation reads
\[
  \bX_1^\top\bX_1\,\bhat_1 + \bX_1^\top\bX_2\,\bhat_2
  = \bX_1^\top\by.
\]
Since $\bX_1$ has full column rank, $\bX_1^\top\bX_1$ is invertible, giving
\begin{equation}\label{eq:fwl-beta1}
  \bhat_1
  = (\bX_1^\top\bX_1)^{-1}\bX_1^\top\by
    - (\bX_1^\top\bX_1)^{-1}\bX_1^\top\bX_2\,\bhat_2.
\end{equation}

\noindent\textbf{Step 2: Substitute into the second block.}
The second block equation reads
\[
  \bX_2^\top\bX_1\,\bhat_1 + \bX_2^\top\bX_2\,\bhat_2
  = \bX_2^\top\by.
\]
Substituting~\eqref{eq:fwl-beta1}:
\begin{align*}
  \bX_2^\top\bX_1
  \bigl[(\bX_1^\top\bX_1)^{-1}\bX_1^\top\by
        &- (\bX_1^\top\bX_1)^{-1}\bX_1^\top\bX_2\,\bhat_2\bigr]
  + \bX_2^\top\bX_2\,\bhat_2
  = \bX_2^\top\by.
\end{align*}
Expanding and rearranging:
\begin{align*}
  &\bX_2^\top\underbrace{\bX_1(\bX_1^\top\bX_1)^{-1}\bX_1^\top}_{\bI - \bM_1}\by
  - \bX_2^\top\underbrace{\bX_1(\bX_1^\top\bX_1)^{-1}\bX_1^\top}_{\bI - \bM_1}\bX_2\,\bhat_2
  + \bX_2^\top\bX_2\,\bhat_2
  = \bX_2^\top\by.
\end{align*}
Since $\bX_1(\bX_1^\top\bX_1)^{-1}\bX_1^\top = \bI - \bM_1$, we can write:
\begin{align*}
  \bX_2^\top(\bI - \bM_1)\by
  &+ \bigl[\bX_2^\top\bX_2
     - \bX_2^\top(\bI - \bM_1)\bX_2\bigr]\bhat_2
  = \bX_2^\top\by.
\end{align*}

\noindent\textbf{Step 3: Simplify.}
Moving $\bX_2^\top(\bI-\bM_1)\by$ to the right-hand side:
\[
  \bX_2^\top\bM_1\bX_2\,\bhat_2
  = \bX_2^\top\by - \bX_2^\top(\bI-\bM_1)\by
  = \bX_2^\top\bM_1\by.
\]
Here we used $\bX_2^\top\bX_2 - \bX_2^\top(\bI-\bM_1)\bX_2 = \bX_2^\top\bM_1\bX_2$
and $\bX_2^\top\by - \bX_2^\top(\bI-\bM_1)\by = \bX_2^\top\bM_1\by$.

Since $\bM_1$ is symmetric and idempotent ($\bM_1^2 = \bM_1$, $\bM_1^\top = \bM_1$),
we have
\[
  \bX_2^\top\bM_1\bX_2 = \bX_2^\top\bM_1^\top\bM_1\bX_2
  = (\bM_1\bX_2)^\top(\bM_1\bX_2),
\]
which is the Gram matrix of the residualised predictors $\bM_1\bX_2$.
Similarly, $\bX_2^\top\bM_1\by = (\bM_1\bX_2)^\top(\bM_1\by)$.

Therefore, provided $\bM_1\bX_2$ has full column rank (i.e., $\bX_2$ is not
in $\Col(\bX_1)$),
\[
  \boxed{\bhat_2
  = \bigl[(\bM_1\bX_2)^\top(\bM_1\bX_2)\bigr]^{-1}(\bM_1\bX_2)^\top(\bM_1\by),}
\]
which is exactly the OLS estimator from regressing the residualised response
$\bM_1\by$ on the residualised predictors $\bM_1\bX_2$.
\end{proof}


% =====================================================================
% ITEM 3: Proper Gauss--Markov Proof
% INSERT: Replace the current "Geometric proof sketch" of Theorem
%         (Gauss--Markov) in Section 4
% =====================================================================

\begin{theorem}[Gauss--Markov]\label{thm:gm-full}
Assume $E[\bm{\varepsilon}] = \bm{0}$ and
$\text{Var}(\bm{\varepsilon}) = \sigma^2\bI_n$
(no assumption on the distribution of $\bm{\varepsilon}$).
Among all linear unbiased estimators of $\bm{c}^\top\bm{\beta}$
(for any fixed $\bm{c}\in\R^p$), the OLS estimator
$\bm{c}^\top\bhat = \bm{c}^\top(\bX^\top\bX)^{-1}\bX^\top\by$
has the smallest variance.
In the vector form: for any linear unbiased estimator
$\btilde = \bm{C}\by$ of $\bm{\beta}$,
\[
  \text{Var}(\btilde) - \text{Var}(\bhat) \;\succeq\; \bm{0}
  \qquad\text{(positive semi-definite).}
\]
\end{theorem}

\begin{proof}
\textbf{Step 1: Set up the alternative estimator.}
Let $\btilde = \bm{C}\by$ be any linear estimator of $\bm{\beta}$, where
$\bm{C}\in\R^{p\times n}$.  Write
\[
  \bm{C} = (\bX^\top\bX)^{-1}\bX^\top + \bD,
\]
where $\bD = \bm{C} - (\bX^\top\bX)^{-1}\bX^\top$ is the ``excess'' beyond
the OLS matrix.

\medskip
\noindent\textbf{Step 2: Derive the constraint from unbiasedness.}
For $\btilde$ to be unbiased for $\bm{\beta}$, we need
$E[\btilde] = \bm{\beta}$ for \emph{every} $\bm{\beta}\in\R^p$.
Since $\by = \bX\bm{\beta} + \bm{\varepsilon}$:
\begin{align*}
  E[\btilde]
  &= \bm{C}\,E[\by]
   = \bm{C}\bX\bm{\beta}
   = \bigl[(\bX^\top\bX)^{-1}\bX^\top + \bD\bigr]\bX\bm{\beta} \\
  &= \bI_p\bm{\beta} + \bD\bX\bm{\beta}
   = (\bI_p + \bD\bX)\bm{\beta}.
\end{align*}
For this to equal $\bm{\beta}$ for every $\bm{\beta}$, we need
\begin{equation}\label{eq:gm-constraint}
  \boxed{\bD\bX = \bm{0}.}
\end{equation}
This is the key constraint: the rows of $\bD$ must be orthogonal to the
columns of $\bX$.  In other words, $\bD$ maps the column space $\Col(\bX)$
to zero.

\medskip
\noindent\textbf{Step 3: Compute the variance and decompose.}
The variance of $\btilde$ is
\begin{align*}
  \text{Var}(\btilde)
  &= \bm{C}\,\text{Var}(\by)\,\bm{C}^\top
   = \sigma^2\bm{C}\bm{C}^\top \\
  &= \sigma^2\bigl[(\bX^\top\bX)^{-1}\bX^\top + \bD\bigr]
              \bigl[\bX(\bX^\top\bX)^{-1} + \bD^\top\bigr].
\end{align*}
Expanding the product:
\begin{align*}
  \bm{C}\bm{C}^\top
  &= (\bX^\top\bX)^{-1}\underbrace{\bX^\top\bX}_{}(\bX^\top\bX)^{-1}
     + (\bX^\top\bX)^{-1}\bX^\top\bD^\top
     + \bD\bX(\bX^\top\bX)^{-1}
     + \bD\bD^\top \\
  &= (\bX^\top\bX)^{-1}
     + (\bX^\top\bX)^{-1}\underbrace{(\bD\bX)^\top}_{\bm{0}}
     + \underbrace{\bD\bX}_{\bm{0}}(\bX^\top\bX)^{-1}
     + \bD\bD^\top,
\end{align*}
where the cross terms vanish by the constraint $\bD\bX = \bm{0}$.
Therefore:
\[
  \text{Var}(\btilde)
  = \sigma^2(\bX^\top\bX)^{-1} + \sigma^2\bD\bD^\top
  = \text{Var}(\bhat) + \sigma^2\bD\bD^\top.
\]

\medskip
\noindent\textbf{Step 4: Conclude.}
The matrix $\bD\bD^\top$ is positive semi-definite (for any vector $\bm{v}$,
$\bm{v}^\top\bD\bD^\top\bm{v} = \norm{\bD^\top\bm{v}}^2 \ge 0$).
Therefore
\[
  \boxed{
  \text{Var}(\btilde) - \text{Var}(\bhat)
  = \sigma^2\bD\bD^\top \succeq \bm{0}.
  }
\]
Equality holds if and only if $\bD = \bm{0}$, i.e., $\bm{C} = (\bX^\top\bX)^{-1}\bX^\top$.
The OLS estimator is the unique Best Linear Unbiased Estimator (BLUE).
\end{proof}

\begin{remark}[Geometric reading of Gauss--Markov]
The constraint $\bD\bX = \bm{0}$ says that the rows of $\bD$ live in
$\Col(\bX)^\perp$.  Any linear unbiased estimator that differs from OLS
must ``reach into'' the orthogonal complement of the column space to
form its estimate.  But $\Col(\bX)^\perp$ is exactly where the noise lives
(the residual space).  The extra term $\sigma^2\bD\bD^\top$ is the
variance penalty for picking up noise from outside the column space.
OLS is optimal precisely because it ignores everything orthogonal to
$\Col(\bX)$.
\end{remark}


% =====================================================================
% ITEM 4: Characterisation Theorem (Symmetric + Idempotent <=> Orthogonal
%         Projection onto Column Space)
% INSERT: In Section 3 (The Hat Matrix and the Right Angle), after the
%         proof of Proposition (Properties of H), before the "Geometric
%         meaning" paragraph. This replaces and subsumes the current
%         informal geometric-meaning paragraph.
% =====================================================================

\begin{theorem}[Characterisation of Orthogonal Projections]\label{thm:proj-char}
Let $\bm{P}\in\R^{n\times n}$.  The following are equivalent:
\begin{enumerate}[nosep,label=(\alph*)]
  \item $\bm{P}$ is the matrix of orthogonal projection onto $\Col(\bm{P})$.
  \item $\bm{P}^2 = \bm{P}$ and $\bm{P}^\top = \bm{P}$.
\end{enumerate}
\end{theorem}

\begin{proof}
\textbf{(a) $\Rightarrow$ (b).}
Let $W = \Col(\bm{P})$ and suppose $\bm{P}$ is the orthogonal projection
onto $W$.  For any $\bm{v}\in\R^n$, $\bm{P}\bm{v}\in W$ and
$\bm{v} - \bm{P}\bm{v}\in W^\perp$.

\emph{Idempotency:}
Since $\bm{P}\bm{v}\in W$, projecting it again gives
$\bm{P}(\bm{P}\bm{v}) = \bm{P}\bm{v}$. Hence $\bm{P}^2 = \bm{P}$.

\emph{Symmetry:}
For any $\bm{u}, \bm{v}\in\R^n$,
\begin{align*}
  \ip{\bm{P}\bm{u}}{\bm{v}}
  &= \ip{\bm{P}\bm{u}}{\bm{P}\bm{v} + (\bm{v}-\bm{P}\bm{v})} \\
  &= \ip{\bm{P}\bm{u}}{\bm{P}\bm{v}} + \ip{\bm{P}\bm{u}}{\bm{v}-\bm{P}\bm{v}} \\
  &= \ip{\bm{P}\bm{u}}{\bm{P}\bm{v}},
\end{align*}
because $\bm{P}\bm{u}\in W$ and $\bm{v}-\bm{P}\bm{v}\in W^\perp$.
By the same argument with roles swapped,
$\ip{\bm{u}}{\bm{P}\bm{v}} = \ip{\bm{P}\bm{u}}{\bm{P}\bm{v}}$.
Therefore $\ip{\bm{P}\bm{u}}{\bm{v}} = \ip{\bm{u}}{\bm{P}\bm{v}}$ for all
$\bm{u},\bm{v}$, which means $\bm{P}^\top = \bm{P}$.

\medskip
\noindent\textbf{(b) $\Rightarrow$ (a).}
Assume $\bm{P}^2 = \bm{P}$ and $\bm{P}^\top = \bm{P}$.
Let $W = \Col(\bm{P})$.
We must show: for every $\bm{v}\in\R^n$,
(i) $\bm{P}\bm{v}\in W$, and
(ii) $\bm{v} - \bm{P}\bm{v} \perp W$.

Part (i) is immediate: $\bm{P}\bm{v}$ is a column of $\bm{P}$ (times scalars),
hence in $\Col(\bm{P}) = W$.

For part (ii), let $\bm{w}\in W$ be arbitrary. Then $\bm{w} = \bm{P}\bm{z}$
for some $\bm{z}$.  Compute:
\[
  \ip{\bm{v} - \bm{P}\bm{v}}{\bm{w}}
  = \ip{\bm{v} - \bm{P}\bm{v}}{\bm{P}\bm{z}}
  = (\bm{v} - \bm{P}\bm{v})^\top\bm{P}\bm{z}
  = \bm{v}^\top\bm{P}\bm{z} - \bm{v}^\top\bm{P}^\top\bm{P}\bm{z}.
\]
Using $\bm{P}^\top = \bm{P}$ and then $\bm{P}^2 = \bm{P}$:
\[
  = \bm{v}^\top\bm{P}\bm{z} - \bm{v}^\top\bm{P}^2\bm{z}
  = \bm{v}^\top\bm{P}\bm{z} - \bm{v}^\top\bm{P}\bm{z}
  = 0.
\]
Since $\bm{w}\in W$ was arbitrary, $(\bm{v}-\bm{P}\bm{v})\perp W$.
\end{proof}

\begin{remark}[Why this theorem matters]
This theorem is the geometric backbone of all projection-based arguments in
linear regression.  It tells us that verifying two algebraic properties
(symmetry and idempotency) is \emph{necessary and sufficient} to conclude
that a matrix is an orthogonal projection.
\begin{itemize}[nosep]
  \item The hat matrix $\bH = \bX(\bX^\top\bX)^{-1}\bX^\top$ satisfies both
        properties (Proposition~\ref{prop:hat}), confirming it is the orthogonal
        projection onto $\Col(\bX)$.
  \item The annihilator $\bM = \bI - \bH$ likewise satisfies both, confirming
        it is the orthogonal projection onto $\Col(\bX)^\perp$.
  \item Idempotency alone (\emph{without} symmetry) gives an \emph{oblique}
        projection: the decomposition $\bm{v} = \bm{P}\bm{v} + (\bm{v}-\bm{P}\bm{v})$
        still holds, but the two components are not perpendicular.
        Symmetry is what makes the projection \emph{orthogonal}---it ensures
        the ``shadow falls straight down'' rather than at an angle.
\end{itemize}
\end{remark}


% =====================================================================
% ITEM 5a: Full Column Rank Assumption Box
% INSERT: At the very beginning of Section 2, immediately after the
%         section heading "Where Coefficients Come From", before
%         \subsection{Setup and notation}
% =====================================================================

\begin{tcolorbox}[
  colback=colspace!8, colframe=colspace!80!black,
  fonttitle=\bfseries, title=Standing Assumption,
  arc=2mm, boxrule=0.8pt
]
\textbf{Full column rank.}
Throughout these notes (unless stated otherwise), we assume that the
$n\times p$ design matrix $\bX$ has \textbf{full column rank}:
$\rank(\bX) = p$, with $n \ge p$.

This is equivalent to any of the following:
\begin{enumerate}[nosep,label=(\roman*)]
  \item The columns of $\bX$ are linearly independent.
  \item The Gram matrix $\bX^\top\bX\in\R^{p\times p}$ is invertible
        (positive definite).
  \item The column space $\Col(\bX)$ has dimension exactly $p$.
  \item No predictor is an exact linear combination of the others.
\end{enumerate}
Without this assumption, $(\bX^\top\bX)^{-1}$ does not exist, and the OLS
coefficient vector $\bhat$ is not uniquely determined (although the
fitted values $\yhat = \bH\by$ remain unique as a projection).
\end{tcolorbox}


% =====================================================================
% ITEM 5b: Fix h_ii >= 1/n (requires intercept)
% INSERT: Replace the current Definition of leverage in Section 5
%         (Leverage and Influence), specifically the claim
%         "1/n <= h_{ii} <= 1"
% =====================================================================

\begin{definition}
The \textbf{leverage} of observation $i$ is $h_{ii}$, the $i$-th diagonal
element of $\bH$. By Proposition~\ref{prop:hat}(iv),
$0 \le h_{ii} \le 1$ and $\sum_i h_{ii} = p$.
The conventional threshold for ``high leverage'' is $h_{ii} > 2p/n$.
\end{definition}

\begin{remark}[The bound $h_{ii}\ge 1/n$ when an intercept is present]\label{rem:hii-intercept}
When the model includes an intercept (i.e., a column of ones $\bm{1}\in\Col(\bX)$),
the leverage satisfies the stronger bound $h_{ii}\ge 1/n$.
To see why: the hat matrix projects onto $\Col(\bX)$, which contains $\bm{1}$.
In particular, $\bH\bm{1} = \bm{1}$, so the $i$-th row of $\bH$ sums to 1:
$\sum_{j=1}^n h_{ij} = 1$.
Since $\bH$ is an orthogonal projection, $h_{ij}^2 \le h_{ii}\cdot h_{jj} \le h_{ii}$
(by Cauchy--Schwarz applied to the rows/columns of $\bH$).  In particular,
\[
  1 = \Bigl(\sum_{j=1}^n h_{ij}\Bigr)^2
  \le n\sum_{j=1}^n h_{ij}^2
  = n\,(\bH^2)_{ii}
  = n\,h_{ii},
\]
where the inequality is Cauchy--Schwarz and the last step uses $\bH^2 = \bH$.
Therefore $h_{ii}\ge 1/n$.

Without an intercept, it is possible to have $h_{ii} = 0$ (for an observation
$\bm{x}_i$ orthogonal to every column of $\bX$), so the general bound is
$0 \le h_{ii} \le 1$.
\end{remark}


% =====================================================================
% ITEM 5c: Fix "symmetry treats all directions equally"
% INSERT: Replace the current paragraph that begins "Symmetry says
%         the projection treats all directions equally..." in Section 3,
%         after the proof of Proposition (Properties of H)
% =====================================================================

% REPLACE the paragraph:
%   "Symmetry says the projection treats all directions equally---it
%    flattens onto the plane without distortion."
% WITH:

\noindent\textbf{Geometric meaning.}
Idempotency says: \emph{the shadow of a shadow is itself.}
If $\yhat$ is already on the column space, projecting it again
changes nothing.
Symmetry says the projection is \emph{orthogonal} (perpendicular), not
oblique.
A non-symmetric idempotent matrix would still decompose
$\by = \bm{P}\by + (\bI - \bm{P})\by$ into a component in $\Col(\bm{P})$
and a remainder, but the two components would not be perpendicular---the
``shadow'' would fall at a slant.
Symmetry of $\bH$ is what guarantees that the residual $\be = (\bI-\bH)\by$
is truly perpendicular to the column space, so that
$\ip{\yhat}{\be} = \by^\top\bH^\top(\bI-\bH)\by = \by^\top(\bH - \bH^2)\by = 0$.
In short: \emph{idempotency makes it a projection; symmetry makes it orthogonal.}


% =====================================================================
% ITEM 5d: Remark on SST = SSR + SSE requiring intercept
% INSERT: After the variance decomposition equation
%         SST = SSR + SSE in Section 4 (Pythagoras Runs a Regression),
%         after "The variance decomposition is the Pythagorean theorem..."
% =====================================================================

\begin{remark}[The intercept is needed for SST $=$ SSR $+$ SSE]\label{rem:intercept-decomp}
The decomposition $\text{SST} = \text{SSR} + \text{SSE}$ requires that the
model includes an intercept (equivalently, that $\bm{1}\in\Col(\bX)$).
Here is why.

The decomposition uses demeaned vectors: $\tilde{\by} = \by - \bar{y}\bm{1}$
and $\tilde{\yhat} = \yhat - \bar{y}\bm{1}$.  The Pythagorean theorem gives
$\norm{\tilde{\by}}^2 = \norm{\tilde{\yhat}}^2 + \norm{\be}^2$
only if $\tilde{\yhat}\perp\be$.

Now $\tilde{\yhat} = \yhat - \bar{y}\bm{1}$.
We know $\yhat\perp\be$ is \emph{not} generally true (it holds only when
$\by\in\Col(\bX)$).
What we need is $(\yhat - \bar{y}\bm{1})\perp\be$, which expands to
$\yhat^\top\be = \bar{y}\bm{1}^\top\be$.

We always have $\yhat^\top\be = 0$ (since $\yhat\in\Col(\bX)$ and
$\be\perp\Col(\bX)$).  So we need $\bar{y}\bm{1}^\top\be = 0$.
If $\bar{y}\ne 0$, this requires $\bm{1}^\top\be = 0$, i.e., that the
residuals sum to zero.

The residuals sum to zero if and only if $\bm{1}\in\Col(\bX)$
(because $\be\perp\Col(\bX)$ and $\bm{1}\in\Col(\bX)$ implies
$\bm{1}^\top\be = 0$).
Without an intercept, residuals need not sum to zero, and the clean
decomposition $\text{SST} = \text{SSR} + \text{SSE}$ can fail.

In models without an intercept, $R^2 = 1 - \text{SSE}/\text{SST}$
can even be negative, which is a symptom of this geometric misalignment.
\end{remark}


% =====================================================================
% ITEM 6: Adjusted R^2 with Geometric Interpretation
% INSERT: After the R^2 monotonicity warning in Section 4, before
%         Section 5 (Peeling Apart the Variables)
% =====================================================================

\subsection{Adjusted $R^2$: paying for dimensions}\label{sec:adj-r2}

The monotonicity of $R^2$ (Proposition~\ref{prop:r2mono}) means that
$R^2$ can never decrease when we add predictors, even useless ones.
This makes raw $R^2$ unreliable for model comparison.
The \textbf{adjusted $R^2$} corrects for this by penalising model complexity.

\begin{definition}[Adjusted $R^2$]
\[
  R^2_{\text{adj}}
  = 1 - \frac{\text{SSE}/(n-p)}{\text{SST}/(n-1)}
  = 1 - \frac{n-1}{n-p}\,(1-R^2).
\]
\end{definition}

\noindent\textbf{Algebraic reading.}
Instead of comparing raw sums of squares (SSE vs SST), adjusted $R^2$
compares \emph{mean squares}: it divides each by its degrees of freedom
($n-p$ for the residuals, $n-1$ for the total).
Adding a useless predictor decreases SSE only by the amount expected from
fitting noise in one extra dimension, while $n-p$ decreases by 1.
The ratio $\text{SSE}/(n-p)$ may actually increase, causing
$R^2_{\text{adj}}$ to decrease.  Thus $R^2_{\text{adj}}$, unlike $R^2$,
can go down when a predictor is added.

\medskip
\noindent\textbf{Geometric reading.}
Think of the data vector $\by$ living in $\R^n$.
The column space $\Col(\bX)$ has dimension $p$, and the residual space
$\Col(\bX)^\perp$ has dimension $n - p$.
When we project $\by$ onto a $p$-dimensional subspace, we are ``using up''
$p$ dimensions to capture signal.  Even pure noise, projected onto a
$p$-dimensional subspace, produces a nonzero shadow---on average, the
squared length of this shadow is $p\sigma^2$.

The residual $\be$ lives in the remaining $n-p$ dimensions.
If the model captures no real signal, we expect
$\norm{\be}^2 \approx (n-p)\sigma^2$ (noise spread over $n-p$ dimensions)
and $\text{SST} \approx (n-1)\sigma^2$ (noise spread over $n-1$ dimensions,
after demeaning).
The ratio $\frac{\text{SSE}/(n-p)}{\text{SST}/(n-1)}$ then hovers around 1,
giving $R^2_{\text{adj}} \approx 0$---correctly indicating no explanatory power.

In short, $R^2_{\text{adj}}$ accounts for the fact that \emph{projecting onto
a subspace costs degrees of freedom}: each dimension consumed by the model
must ``earn its keep'' by reducing the residual more than random chance would.


% =====================================================================
% ITEM 7: Degrees of Freedom Section
% INSERT: As a new subsection in Section 4 (Pythagoras Runs a Regression),
%         after the adjusted R^2 subsection and before Section 5 (FWL).
%         Alternatively, as a standalone subsection in Section 3 after the
%         annihilator matrix subsection.
% =====================================================================

\subsection{Degrees of freedom: the dimension of the residual space}\label{sec:df}

The phrase ``degrees of freedom'' appears throughout statistics, often without
geometric motivation.  The projection framework makes it transparent.

\begin{definition}[Degrees of freedom of the residual]
In a linear model with $p$ parameters fitted to $n$ observations, the
\textbf{residual degrees of freedom} is $n - p$.
\end{definition}

\noindent\textbf{Geometric interpretation.}
The data vector $\by$ lives in $\R^n$.
The projection $\yhat = \bH\by$ lives in $\Col(\bX)$, which has dimension $p$.
The residual $\be = \bM\by$ lives in $\Col(\bX)^\perp$, which has dimension
\[
  \dim\bigl(\Col(\bX)^\perp\bigr) = n - \dim\bigl(\Col(\bX)\bigr) = n - p.
\]
This is confirmed algebraically: $\rank(\bM) = \tr(\bM) = \tr(\bI_n - \bH) = n - p$.

The residual vector $\be\in\R^n$ has $n$ components, but it satisfies $p$
linear constraints ($\bX^\top\be = \bm{0}$, i.e., $\be\perp\Col(\bX)$).
So $\be$ is free to vary in only $n - p$ dimensions.
That is the number of ``degrees of freedom'' left after fitting the model.

\begin{remark}[Why divide by $n-p$, not $n$]
The unbiased estimator of $\sigma^2$ is
\[
  s^2 = \frac{\norm{\be}^2}{n - p} = \frac{\text{SSE}}{n-p},
\]
not $\norm{\be}^2/n$.  The geometric reason: $\be = \bM\bm{\varepsilon}$
(since $\bM\bX = \bm{0}$), and the expected squared length of
$\bm{\varepsilon}$ projected onto a $(n-p)$-dimensional subspace is
\[
  E[\norm{\bM\bm{\varepsilon}}^2]
  = E[\bm{\varepsilon}^\top\bM\bm{\varepsilon}]
  = \sigma^2\tr(\bM)
  = (n-p)\sigma^2.
\]
(Here we used $E[\bm{\varepsilon}\bm{\varepsilon}^\top] = \sigma^2\bI$
and the trace-expectation identity.)
Dividing by $n-p$ exactly compensates for the dimensionality of the
space into which the noise is projected, yielding $E[s^2] = \sigma^2$.

Dividing by $n$ would systematically underestimate $\sigma^2$, because
$p$ dimensions of the noise have been ``absorbed'' by the projection
onto $\Col(\bX)$---the hat matrix captures $p$ dimensions worth of noise
as spurious fit.
\end{remark}

\begin{remark}[Degrees-of-freedom bookkeeping]
The total degrees of freedom decompose just like the sums of squares:
\[
  \underbrace{n-1}_{\text{total d.f.}}
  \;=\;
  \underbrace{p-1}_{\text{model d.f.}}
  \;+\;
  \underbrace{n-p}_{\text{residual d.f.}},
\]
where we subtract 1 from each side because the intercept uses one dimension
to capture the mean (assuming the model includes an intercept).
This mirrors the Pythagorean decomposition: each sum of squares is
``spread'' over its corresponding number of dimensions, and dividing gives
the mean square:
\[
  \text{MSR} = \frac{\text{SSR}}{p-1},
  \qquad
  \text{MSE} = \frac{\text{SSE}}{n-p}.
\]
The $F$-statistic (Section~\ref{sec:ftest}) is the ratio $\text{MSR}/\text{MSE}$:
signal per dimension versus noise per dimension.
\end{remark}
